\cdsection{The Effects of Deforestation}

\cdstep{A}

Every year it is estimated that roughly 5.2 million hectares (52,000 km2
of the forest is lost worldwide. That is a net figure, meaning it
represents the area of the forest not replaced. T put this size in
context, that is an area of land the size of Croatia lost every single
year. There is a wide range of negative effects from deforestation that
range from the smallest biological processes right up to the health of
our planet as a whole. On a human level, millions of lives are affected
every year by flooding and landslides that often result from
deforestation.

\cdstep{B}

There are 5 million people living in areas deemed at risk of flooding in
England and Wales. Global warming, in part, worsened by deforestation,
is responsible for higher rainfalls in Britain in recent decades.
Although it can be argued that demand for cheap housing has meant more
houses are being built in at-risk areas, the extent of the flooding is
increasing. The presence of forests and trees along streams and rivers
acts like a net. The trees catch and store water, but also hold the soil
together, preventing erosion. By removing the trees, the land is more
easily eroded increasing the risk of landslides and also, after
precipitation, less water is intercepted when trees are absent and so
more enters rivers, increasing the risk of flooding.

\cdstep{C}

It is well documented that forests are essential to the atmospheric
balance of our planet, and therefore our own wellbeing too. Scientists
agree unequivocally that global warming is a real and serious threat to
our planet. Deforestation releases 15\% of all greenhouse gas emissions.
One-third of the carbon dioxide emissions created by human activity
comes from deforestation around the globe.

\cdstep{D}

In his book Collapse, about the disappearance of various ancient
civilisations, writer Jared Diamond theorises about the decline of the
natives of Easter Island. European missionaries first arrived on the
island in 1722. Research suggested that the island, whose population was
in the region of two to three thousand at the time, had once been much
higher at fifteen thousand people. This small native population survived
on the island despite there being no trees at all. Archaeological digs
uncovered evidence of trees once flourishing on the island. The
uncontrolled deforestation not only led to the eradication of all such
natural resources from the island but also greatly impacted the number
of people the island could sustain. This underlines the importance of
forest management, not only for useful building materials but also for
food as well.

\cdstep{E}

Forestry management is important to make sure that stocks are not
depleted and that whatever is cut down is replaced. Without sustainable
development of forests, the levels of deforestation are only going to
worsen as the global population continues to rise, creating a higher
demand for the products of forests. Just as important though is consumer
awareness. Simple changes in consumer activity can make a huge
difference. These changes in behaviour include, but are not limited to,
recycling all recyclable material; buying recycled products and looking
for the FSC sustainably sourced forest products logo on any wood or
paper products.

\cdstep{F}

Japan is often used as a model of exemplary forest management. During
the Edo period between 1603 and 1868 drastic action was taken to reverse
the country's serious exploitative deforestation problem. Whilst the
solution was quite complex, one key aspect of its success was the
encouragement of cooperation between villagers. This process of
collaboration and re-education of the population saved Japan's forests.
According to the World Bank, 68.5\% of Japanese land area is covered by
forest, making it one of the best performing economically developed
nations in this regard.

\cdstep{G}

There is, of course, a negative impact of Japan's forest management.
There is still a high demand for wood products in the country, and the
majority of these resources are simply imported from other, poorer
nations. Indonesia is a prime example of a country that has lost large
swaths of its forest cover due to foreign demand from countries like
Japan. This is in addition to other issues such as poor domestic forest
management, weaker laws and local corruption. Located around the
Equator, Indonesia has an ideal climate for the rainforest. Sadly much
of this natural resource is lost every year. Forest cover is now down to
less than 51\% from 65.4\% in 1990. This alone is proof that more needs
to be done globally to manage forests.

\cdstep{H}

China is leading the way in recent years for replenishing their forests.
The Chinese government began the Three-North Shelter Forest Program in
1978, with aims to complete the planting of a green wall, measuring
2,800 miles in length by its completion in 2050. Of course, this program
is in many ways forced by nature itself; the expansion of the Gobi
Desert threatened to destroy thousands of square miles of grassland
annually through desertification. This is a process often exacerbated by
deforestation in the first place, and so represents an attempt to buck
the trend. Forested land in China rose from 17\% to 22\% from 1990 to
2015 making China one of the few developing nations to reverse the
negative trend.

\cdsection{Glossary}

\begin{itemize}
\tightlist
\item
  \textbf{exemplary:} serving as a perfect example
\item
  \textbf{exacerbate:} make worse
\end{itemize}

\cdsection{Questions 14-20}

The reading passage below has eight paragraphs, A-H.

\begin{cdnote}

Choose the correct heading for each paragraph from the list of headings
below.

\end{cdnote}

\begin{cdnote}

Write the correct number, i-xi, in boxes 14-20 on your answer sheet.

\end{cdnote}

\cdstep{Example Answer}

\begin{itemize}
\tightlist
\item
  \textbf{Paragraph A:} vi
\end{itemize}

\begingroup
\sloppy
\renewcommand{\tabularxcolumn}[1]{>{\RaggedRight\arraybackslash}m{#1}}
\setlength{\tabcolsep}{6pt}
\renewcommand{\arraystretch}{1.15}
\arrayrulecolor{black!25}
\begin{cdtablebox}
\begin{tabularx}{\linewidth}{|M{0.18\linewidth}|Y|}
\hline
i & Atmospheric impacts \\ \hline
ii & Ideal forestry management example \\ \hline
iii & No trees, less people \\ \hline
iv & Good uses for wood \\ \hline
v & Looking after the forests \\ \hline
vi & Numbers of lost trees \\ \hline
vii & Wasted water \\ \hline
viii & Replanting forests \\ \hline
ix & Happy trees \\ \hline
x & Flood risks \\ \hline
xi & Poorer nations at higher risk \\ \hline
\end{tabularx}
\end{cdtablebox}
\endgroup

\begin{cdquestions}

\textbf{14.} Paragraph B\\
\textbf{15.} Paragraph C\\
\textbf{16.} Paragraph D\\
\textbf{17.} Paragraph E\\
\textbf{18.} Paragraph F\\
\textbf{19.} Paragraph G\\
\textbf{20.} Paragraph H

\end{cdquestions}

\cdsection{Questions 21-26}

\begin{cdnote}

Complete the summary below.

\end{cdnote}

\begin{cdnote}

Choose NO MORE THAN TWO WORDS AND/OR A NUMBER from the passage for each
answer.

\end{cdnote}

\begin{cdnote}

Write your answers in boxes 21-26 on your answer sheet.

\end{cdnote}

The effects of deforestation are widespread and various. Some examples
include flooding at a local scale to the wider effects of global warming
on a worldwide scale. In Britain, for example, 21
\ldots\ldots\ldots\ldots\ldots\ldots\ldots. people live in areas at risk
of flooding. This risk is increased by deforestation. Trees catch and 22
\ldots\ldots\ldots\ldots\ldots\ldots\ldots{} water lowering the chance
of flooding. By removing trees land erosion is also higher, increasing
the chance of 23 \ldots\ldots\ldots\ldots\ldots\ldots\ldots{}
Deforestation also affects global warming by contributing 15\% of the 24
\ldots\ldots\ldots\ldots\ldots\ldots\ldots. of greenhouse gasses. To
make sure that the cutting down of trees is done in a sustainable way,
good forestry 25 \ldots\ldots\ldots\ldots\ldots\ldots\ldots. is
important. In most countries, more trees are cut down every year than
planted. One country that is reversing this trend is China, making it
one of the few nations to 26
\ldots\ldots\ldots\ldots\ldots\ldots\ldots.. the more common negative
trend.

\textbf{\cdblue{Hướng dẫn giải}}

\cdsection{Phần I: Câu hỏi 14-20 (Matching Headings)}

\textbf{Mục tiêu:} Nắm bắt ý chính của đoạn văn và tìm tiêu đề tương ứng
trong danh sách cho sẵn.

\cdstep{1.\allowbreak{} Chiến lược tiếp cận}

\begin{itemize}
\tightlist
\item
  \textbf{Bước 1:} Giải mã Danh sách Tiêu đề (List of Headings)

  \begin{itemize}
  \tightlist
  \item
    Đọc kỹ danh sách các tiêu đề từ (i) đến (xi).
  \item
    Gạch chân từ khóa và tóm lược nội dung chính:
  \end{itemize}
\end{itemize}

\begingroup
\sloppy
\renewcommand{\tabularxcolumn}[1]{>{\RaggedRight\arraybackslash}m{#1}}
\setlength{\tabcolsep}{6pt}
\renewcommand{\arraystretch}{1.15}
\arrayrulecolor{black!25}
\begin{cdtablebox}
\begin{tabularx}{\linewidth}{|M{0.18\linewidth}|Y|}
\hline
i & Atmospheric impacts (Tác động đến khí quyển).\allowbreak{} \\ \hline
ii & Ideal forestry management example (Ví dụ quản lý rừng lý tưởng).\allowbreak{} \\ \hline
iii & No trees, less people (Không có cây, ít người hơn).\allowbreak{} \\ \hline
iv & Good uses for wood (Công dụng tốt của gỗ).\allowbreak{} \\ \hline
v & Looking after the forests (Chăm sóc rừng).\allowbreak{} \\ \hline
vi & Numbers of lost trees (Số lượng cây bị mất).\allowbreak{} \\ \hline
vii & Wasted water (Nước bị lãng phí).\allowbreak{} \\ \hline
viii & Replanting forests (Trồng lại rừng).\allowbreak{} \\ \hline
ix & Happy trees (Cây hạnh phúc).\allowbreak{} \\ \hline
x & Flood risks (Nguy cơ lũ lụt).\allowbreak{} \\ \hline
xi & Poorer nations at higher risk (Các quốc gia nghèo hơn chịu rủi ro cao hơn).\allowbreak{} \\ \hline
\end{tabularx}
\end{cdtablebox}
\endgroup

\begin{itemize}
\tightlist
\item
  \textbf{Bước 2:} Đọc lướt (Skimming) và Đối chiếu

  \begin{itemize}
  \tightlist
  \item
    Đọc lướt từng đoạn văn để tìm câu chủ đề hoặc ý bao trùm.
  \item
    \textbf{Lưu ý:} Tìm sự tương đồng về ý nghĩa giữa tiêu đề và nội
    dung đoạn văn.
  \end{itemize}
\end{itemize}

\cdstep{2.\allowbreak{} Phân tích chi tiết từng đoạn văn}

\begin{itemize}
\tightlist
\item
  \textbf{Câu 14:} Paragraph B

  \begin{itemize}
  \tightlist
  \item
    \textbf{Nội dung chính:} Đoạn văn mở đầu bằng việc nhắc đến ``5
    million people living in areas deemed at risk of flooding''. Tiếp
    theo giải thích việc chặt cây làm tăng nguy cơ này: ``increasing the
    risk of flooding''.
  \item
    \textbf{Kết luận:} Nội dung tập trung hoàn toàn vào rủi ro ngập lụt,
    khớp với tiêu đề (x) Flood risks.
  \item
    \textbf{Đáp án:} x
  \end{itemize}
\item
  \textbf{Câu 15:} Paragraph C

  \begin{itemize}
  \tightlist
  \item
    \textbf{Nội dung chính:} Đoạn văn đề cập đến ``atmospheric balance''
    (cân bằng khí quyển), ``global warming'' (nóng lên toàn cầu), và
    ``greenhouse gas emissions'' (khí thải nhà kính).
  \item
    \textbf{Kết luận:} Các từ khóa này đều liên quan đến tác động lên
    không khí/khí quyển, khớp với tiêu đề (i) Atmospheric impacts.
  \item
    \textbf{Đáp án:} i
  \end{itemize}
\item
  \textbf{Câu 16:} Paragraph D

  \begin{itemize}
  \tightlist
  \item
    \textbf{Nội dung chính:} Đoạn văn kể về sự suy tàn của người bản địa
    đảo Easter. Tác giả viết: ``The uncontrolled deforestation\ldots{}
    greatly impacted the number of people the island could sustain.''
    (Nạn phá rừng không kiểm soát\ldots{} ảnh hưởng lớn đến số lượng
    người mà hòn đảo có thể nuôi sống).
  \item
    \textbf{Kết luận:} Ý chính là việc mất rừng dẫn đến sự suy giảm dân
    số, khớp với tiêu đề (iii) No trees, less people.
  \item
    \textbf{Đáp án:} iii
  \end{itemize}
\item
  \textbf{Câu 17:} Paragraph E

  \begin{itemize}
  \tightlist
  \item
    \textbf{Nội dung chính:} Đoạn văn bàn về ``Forestry management''
    (quản lý lâm nghiệp) để đảm bảo bền vững. Nó cũng nhắc đến
    ``consumer awareness'' (nhận thức người tiêu dùng) và các hành động
    bảo vệ như ``recycling'' (tái chế), ``sustainably sourced'' (nguồn
    gốc bền vững).
  \item
    \textbf{Kết luận:} Tất cả các hành động này đều nhằm mục đích bảo vệ
    và chăm sóc rừng, khớp với tiêu đề (v) Looking after the forests.
  \item
    \textbf{Đáp án:} v
  \end{itemize}
\item
  \textbf{Câu 18:} Paragraph F

  \begin{itemize}
  \tightlist
  \item
    \textbf{Nội dung chính:} Câu đầu tiên khẳng định: ``Japan is often
    used as a model of exemplary forest management'' (Nhật Bản thường
    được dùng làm mô hình quản lý rừng kiểu mẫu).
  \item
    \textbf{Kết luận:} Từ ``exemplary'' (kiểu mẫu) đồng nghĩa với
    ``ideal'' (lý tưởng), khớp với tiêu đề (ii) Ideal forestry
    management example.
  \item
    \textbf{Đáp án:} ii
  \end{itemize}
\item
  \textbf{Câu 19:} Paragraph G

  \begin{itemize}
  \tightlist
  \item
    \textbf{Nội dung chính:} Đoạn văn nói về tác động tiêu cực của việc
    Nhật Bản nhập khẩu gỗ từ ``poorer nations'' (các quốc gia nghèo hơn)
    như Indonesia. Indonesia bị mất rừng do nhu cầu nước ngoài và tham
    nhũng.
  \item
    \textbf{Kết luận:} Nội dung nhấn mạnh việc các nước nghèo đang chịu
    rủi ro mất rừng cao hơn để phục vụ nước giàu, khớp với tiêu đề (xi)
    Poorer nations at higher risk.
  \item
    \textbf{Đáp án:} xi
  \end{itemize}
\item
  \textbf{Câu 20:} Paragraph H

  \begin{itemize}
  \tightlist
  \item
    \textbf{Nội dung chính:} Đoạn văn nói về Trung Quốc đang dẫn đầu
    trong việc ``replenishing their forests'' (làm đầy lại rừng của họ)
    thông qua chương trình trồng rừng ``Three-North Shelter Forest
    Program''.
  \item
    \textbf{Kết luận:} Hành động trồng lại cây và khôi phục rừng khớp
    với tiêu đề (viii) Replanting forests.
  \item
    \textbf{Đáp án:} viii
  \end{itemize}
\end{itemize}

\cdsection{Phần II: Câu hỏi 21-26 (Summary Completion)}

\begin{cdnote}

Yêu cầu: Hoàn thành bản tóm tắt về tác động của phá rừng với giới hạn
KHÔNG QUÁ 2 TỪ VÀ/HOẶC MỘT SỐ.

\end{cdnote}

\cdstep{1.\allowbreak{} Phân tích Từ khóa và Dự đoán}

Chúng ta cần xác định loại từ cần điền dựa trên ngữ pháp và tìm từ khóa
trong bài đọc để định vị thông tin.

\begin{itemize}
\tightlist
\item
  \textbf{(21)} In Britain, for example, \_\_\_\_\_\_\_ people live in
  areas at risk of flooding.

  \begin{itemize}
  \tightlist
  \item
    \textbf{Từ khóa:} ``Britain'', ``people live'', ``risk of
    flooding''.
  \item
    \textbf{Dự đoán:} Cần một con số chỉ số lượng người.
  \item
    \textbf{Định vị:} Đoạn B, câu đầu tiên: ``There are 5 million people
    living in areas deemed at risk of flooding in England and Wales.''
    (Anh và xứ Wales là đại diện cho Britain).
  \end{itemize}
\item
  \textbf{(22)} Trees catch and \_\_\_\_\_\_\_ water lowering the chance
  of flooding.

  \begin{itemize}
  \tightlist
  \item
    \textbf{Từ khóa:} ``Trees catch'', ``water'', ``lowering chance of
    flooding''.
  \item
    \textbf{Dự đoán:} Cần một động từ mô tả hành động của cây với nước.
  \item
    \textbf{Định vị:} Đoạn B, câu: ``The trees catch and store
    water\ldots{}''
  \end{itemize}
\item
  \textbf{(23)} By removing trees land erosion is also higher,
  increasing the chance of \_\_\_\_\_\_\_.

  \begin{itemize}
  \tightlist
  \item
    \textbf{Từ khóa:} ``removing trees'', ``erosion'', ``increasing
    chance of''.
  \item
    \textbf{Dự đoán:} Cần một danh từ chỉ hậu quả của xói mòn (erosion).
  \item
    \textbf{Định vị:} Đoạn B, câu: ``\ldots increasing the risk of
    landslides\ldots{}''
  \end{itemize}
\item
  \textbf{(24)} Deforestation also affects global warming by
  contributing 15\% of the \_\_\_\_\_\_\_ of greenhouse gasses.

  \begin{itemize}
  \tightlist
  \item
    \textbf{Từ khóa:} ``contributing 15\%'', ``greenhouse gasses''.
  \item
    \textbf{Dự đoán:} Cần một danh từ đi sau ``15\% of the\ldots{}''.
  \item
    \textbf{Định vị:} Đoạn C, câu: ``Deforestation releases 15\% of all
    greenhouse gas emissions.''
  \end{itemize}
\item
  \textbf{(25)} To make sure that the cutting down of trees is done in a
  sustainable way, good forestry \_\_\_\_\_\_\_ is important.

  \begin{itemize}
  \tightlist
  \item
    \textbf{Từ khóa:} ``sustainable way'', ``good forestry'',
    ``important''.
  \item
    \textbf{Dự đoán:} Cần một danh từ ghép với ``forestry''.
  \item
    \textbf{Định vị:} Đoạn E, câu đầu tiên: ``Forestry management is
    important to make sure that stocks are not depleted\ldots{}''
  \end{itemize}
\item
  \textbf{(26)} One country that is reversing this trend is China,
  making it one of the few nations to \_\_\_\_\_\_\_ the more common
  negative trend.

  \begin{itemize}
  \tightlist
  \item
    \textbf{Từ khóa:} ``China'', ``reversing this trend'', ``few nations
    to''.
  \item
    \textbf{Dự đoán:} Cần một động từ chỉ hành động ngược lại với xu
    hướng tiêu cực.
  \item
    \textbf{Định vị:} Đoạn H, câu cuối cùng: ``\ldots making China one
    of the few developing nations to reverse the negative trend.''
  \end{itemize}
\end{itemize}

\cdstep{2.\allowbreak{} Đối chiếu và Kết luận Đáp án}

\begingroup
\sloppy
\renewcommand{\tabularxcolumn}[1]{>{\RaggedRight\arraybackslash}m{#1}}
\setlength{\tabcolsep}{6pt}
\renewcommand{\arraystretch}{1.15}
\arrayrulecolor{black!25}
\setlength{\tabcolsep}{4pt}
\setlength{\LTleft}{0pt}
\setlength{\LTright}{0pt}
\setlength{\LTpre}{0pt}
\setlength{\LTpost}{0pt}
\begin{longtable}{|M{0.10\linewidth}|M{0.66\linewidth}|M{0.14\linewidth}|}
\hline
\rowcolor{topicblue} \cdtableheader{Câu} & \cdtableheader{Phân tích và Dẫn chứng} & \cdtableheader{Đáp án} \\ \hline
\endfirsthead
\hline
\rowcolor{topicblue} \cdtableheader{Câu} & \cdtableheader{Phân tích và Dẫn chứng} & \cdtableheader{Đáp án} \\ \hline
\endhead
\hline
\endfoot
\hline
\endlastfoot
21 & Tìm thấy số lượng người sống trong vùng nguy cơ lũ lụt ở Anh tại đoạn B là 5 triệu.\allowbreak{} & 5 million \\ \hline
22 & Đoạn B nói cây “catch and store water” (giữ và lưu trữ nước).\allowbreak{} & store \\ \hline
23 & Đoạn B chỉ ra việc xói mòn đất làm tăng nguy cơ sạt lở đất (landslides).\allowbreak{} & landslides \\ \hline
24 & Đoạn C nêu rõ phá rừng thải ra 15\% lượng khí thải (emissions) nhà kính.\allowbreak{} & emissions \\ \hline
25 & Đoạn E nhấn mạnh tầm quan trọng của quản lý (management) lâm nghiệp.\allowbreak{} & management \\ \hline
26 & Đoạn H nói Trung Quốc là một trong số ít quốc gia đảo ngược (reverse) xu hướng tiêu cực.\allowbreak{} & reverse \\ \hline
\end{longtable}
\endgroup

\cdsection{Bảng đáp án tổng hợp}

\begingroup
\sloppy
\renewcommand{\tabularxcolumn}[1]{>{\RaggedRight\arraybackslash}m{#1}}
\setlength{\tabcolsep}{6pt}
\renewcommand{\arraystretch}{1.15}
\arrayrulecolor{black!25}
\setlength{\tabcolsep}{4pt}
\setlength{\LTleft}{0pt}
\setlength{\LTright}{0pt}
\setlength{\LTpre}{0pt}
\setlength{\LTpost}{0pt}
\begin{longtable}{|M{0.17\linewidth}|M{0.73\linewidth}|}
\hline
\rowcolor{topicblue} \cdtableheader{Câu} & \cdtableheader{Đáp án} \\ \hline
\endfirsthead
\hline
\rowcolor{topicblue} \cdtableheader{Câu} & \cdtableheader{Đáp án} \\ \hline
\endhead
\hline
\endfoot
\hline
\endlastfoot
14 & x \\ \hline
15 & i \\ \hline
16 & iii \\ \hline
17 & v \\ \hline
18 & ii \\ \hline
19 & xi \\ \hline
20 & viii \\ \hline
21 & 5 million \\ \hline
22 & store \\ \hline
23 & landslides \\ \hline
24 & emissions \\ \hline
25 & management \\ \hline
26 & reverse \\ \hline
\end{longtable}
\endgroup

\cdsection{Tổng hợp từ vựng}

\begingroup
\sloppy
\renewcommand{\tabularxcolumn}[1]{>{\RaggedRight\arraybackslash}m{#1}}
\setlength{\tabcolsep}{6pt}
\renewcommand{\arraystretch}{1.15}
\arrayrulecolor{black!25}
\footnotesize
\setlength{\tabcolsep}{4pt}
\renewcommand{\arraystretch}{1.10}
\setlength{\LTpre}{0pt}
\setlength{\LTpost}{0pt}
\setlength{\tabcolsep}{4pt}
\setlength{\LTleft}{0pt}
\setlength{\LTright}{0pt}
\setlength{\LTpre}{0pt}
\setlength{\LTpost}{0pt}
\begin{longtable}{|M{0.22\linewidth}|M{0.13\linewidth}|M{0.38\linewidth}|M{0.20\linewidth}|}
\hline
\rowcolor{topicblue} \cdtableheader{Từ vựng (Từ loại)} & \cdtableheader{Phiên âm} & \cdtableheader{Nghĩa tiếng Anh (Giải thích chi tiết nghĩa tiếng Việt)} & \cdtableheader{Ví dụ minh họa} \\ \hline
\endfirsthead
\hline
\rowcolor{topicblue} \cdtableheader{Từ vựng (Từ loại)} & \cdtableheader{Phiên âm} & \cdtableheader{Nghĩa tiếng Anh (Giải thích chi tiết nghĩa tiếng Việt)} & \cdtableheader{Ví dụ minh họa} \\ \hline
\endhead
\hline
\endfoot
\hline
\endlastfoot
deforestation (n) & \cdipa{/\allowbreak{}ˌdiːˌfɒr.\allowbreak{}ɪˈsteɪ.\allowbreak{}ʃən/\allowbreak{}} & the cutting down of trees in a large area; the destruction of forests by people (sự phá rừng, sự chặt phá rừng) & Every year it is estimated that roughly 5.\allowbreak{}2 million hectares of the forest is lost worldwide due to deforestation.\allowbreak{} \\ \hline
erosion (n) & \cdipa{/\allowbreak{}ɪˈrəʊ.\allowbreak{}ʒən/\allowbreak{}} & the process by which the surface of the earth is worn away by the action of water, glaciers, winds, waves, etc.\allowbreak{} (sự xói mòn đất) & The trees hold the soil together, preventing erosion.\allowbreak{} \\ \hline
precipitation (n) & \cdipa{/\allowbreak{}prɪˌsɪp.\allowbreak{}ɪˈteɪ.\allowbreak{}ʃən/\allowbreak{}} & rain, snow, sleet, or hail that falls to the ground (lượng mưa, lượng nước rơi xuống) & After precipitation, less water is intercepted when trees are absent.\allowbreak{} \\ \hline
unequivocally (adv) & \cdipa{/\allowbreak{}ˌʌn.\allowbreak{}ɪˈkwɪv.\allowbreak{}ə.\allowbreak{}kəl.\allowbreak{}i/\allowbreak{}} & in a way that is clear and has only one possible meaning; without any doubt (một cách rõ ràng, không thể chối cãi) & Scientists agree unequivocally that global warming is a real and serious threat.\allowbreak{} \\ \hline
emissions (n) & \cdipa{/\allowbreak{}ɪˈmɪʃ.\allowbreak{}ənz/\allowbreak{}} & the production and discharge of something, especially gas or radiation (khí thải, sự phát thải) & Deforestation releases 15\% of all greenhouse gas emissions.\allowbreak{} \\ \hline
decimate (v) & \cdipa{/\allowbreak{}ˈdes.\allowbreak{}ɪ.\allowbreak{}meɪt/\allowbreak{}} & to kill a large number of something, or to reduce something severely (tàn phá, giết hại một lượng lớn) & This epidemic decimated the city until the first frost came.\allowbreak{} \\ \hline
inclement (adj) & \cdipa{/\allowbreak{}ɪnˈklem.\allowbreak{}ənt/\allowbreak{}} & (of weather) unpleasant, cold, and often raining (khắc nghiệt, không thuận lợi - về thời tiết) & The inclement weather froze out the insects.\allowbreak{} \\ \hline
eradication (n) & \cdipa{/\allowbreak{}ɪˌræd.\allowbreak{}ɪˈkeɪ.\allowbreak{}ʃən/\allowbreak{}} & the complete destruction or removal of something (sự xóa bỏ hoàn toàn, sự tiêu diệt) & The uncontrolled deforestation led to the eradication of all such natural resources from the island.\allowbreak{} \\ \hline
sustain (v) & \cdipa{/\allowbreak{}səˈsteɪn/\allowbreak{}} & to cause or allow something to continue for a period of time; to support or keep something alive (duy trì, nuôi sống) & This greatly impacted the number of people the island could sustain.\allowbreak{} \\ \hline
deplete (v) & \cdipa{/\allowbreak{}dɪˈpliːt/\allowbreak{}} & to reduce something in size or amount, especially supplies of energy, money, etc.\allowbreak{} (làm cạn kiệt, làm suy giảm) & Forestry management is important to make sure that stocks are not depleted.\allowbreak{} \\ \hline
exemplary (adj) & \cdipa{/\allowbreak{}ɪɡˈzem.\allowbreak{}plər.\allowbreak{}i/\allowbreak{}} & serving as a desirable model; very good and suitable to be copied (kiểu mẫu, gương mẫu) & Japan is often used as a model of exemplary forest management.\allowbreak{} \\ \hline
exploitative (adj) & \cdipa{/\allowbreak{}ɪkˈsplɔɪ.\allowbreak{}tə.\allowbreak{}tɪv/\allowbreak{}} & unfairly using someone or something for your own advantage (mang tính bóc lột, khai thác quá mức) & Drastic action was taken to reverse the country’s serious exploitative deforestation problem.\allowbreak{} \\ \hline
collaboration (n) & \cdipa{/\allowbreak{}kəˌlæb.\allowbreak{}əˈreɪ.\allowbreak{}ʃən/\allowbreak{}} & the act of working together with someone to create or produce something (sự hợp tác, sự cộng tác) & One key aspect of its success was the encouragement of cooperation between villagers through collaboration.\allowbreak{} \\ \hline
impoverished (adj) & \cdipa{/\allowbreak{}ɪmˈpɒv.\allowbreak{}ər.\allowbreak{}ɪʃt/\allowbreak{}} & very poor; made poor (nghèo khó, bị làm nghèo đi) & Countries with fewer than five frosty days are impoverished.\allowbreak{} \\ \hline
exacerbate (v) & \cdipa{/\allowbreak{}ɪɡˈzæs.\allowbreak{}ə.\allowbreak{}beɪt/\allowbreak{}} & to make something that is already bad even worse (làm trầm trọng thêm, làm tồi tệ hơn) & This is a process often exacerbated by deforestation in the first place.\allowbreak{} \\ \hline
desertification (n) & \cdipa{/\allowbreak{}dɪˌzɜː.\allowbreak{}tɪ.\allowbreak{}fɪˈkeɪ.\allowbreak{}ʃən/\allowbreak{}} & the process by which land changes into desert (sự sa mạc hóa) & The expansion of the Gobi Desert threatened to destroy thousands of square miles of grassland annually through desertification.\allowbreak{} \\ \hline
replenish (v) & \cdipa{/\allowbreak{}rɪˈplen.\allowbreak{}ɪʃ/\allowbreak{}} & to fill something up again by replacing what has been used (bổ sung lại, làm đầy trở lại) & China is leading the way in recent years for replenishing their forests.\allowbreak{} \\ \hline
\end{longtable}
\endgroup
